\section{Intrinsic gravitational constant and tensor representation}
\label{rm:ch:gravity}

\subsection{Constitutive chain of the gravitational sector: chronological scale \texorpdfstring{\(108\)}{108}, tensor reduction, holonomy and metrological evaluation of \texorpdfstring{\(G\)}{G}}

Gravity is organized here into four strata that must not be merged:

\begin{equation}
\label{rm:eq:gravity-causal-chain}
\begin{aligned}
\text{CT108 chronological structure and action lines}
&\longrightarrow \text{periodicity selector }\theta_{108}^{\rm circ}
\longrightarrow G_{108},\\
\text{icosahedral incidence}
&\longrightarrow V_{12}\xrightarrow{P_5}V_5
\xrightarrow{\Gamma_5}\operatorname{STF}_3
\xrightarrow{\Lambda(n)}\mathcal H_{\rm TT}(n),\\
\text{enriched holonomic route}
&\dashrightarrow G_{\rm hol},\\
\text{metrological chart}
&\longrightarrow 6.67430\times10^{-11}\,
\mathrm{m^3\,kg^{-1}\,s^{-2}}.
\end{aligned}
\end{equation}

The first line states a condition of periodicity that is intrinsic and exact.
The second constructs the carrier and its trace-free transverse reduction. The third preserves an
explicit construction boundary. The fourth is an evaluation in a
declared decimal chart: it identifies the magnitude but does not select the gravitational
character.

\subsection{Gravitational magnitude line and chronological scale \texorpdfstring{\(108\)}{108}}

Let \(c_{\rm int}>0\), \(t_0>0\),
\(\ell_0=c_{\rm int}t_0\) and a positive section
\(\hbar_a\) of the action line, with
\(a\in\{\mathrm{pre},\mathrm{ret}\}\). The dimensional transducer is

\begin{equation}
\label{rm:eq:gravity-transducer}
G_a(\theta)
=\theta\frac{c_{\rm int}^{5}t_0^2}{\hbar_a}
=\theta\frac{c_{\rm int}^{3}\ell_0^2}{\hbar_a},
\qquad \theta>0.
\end{equation}

The equality of the two expressions uses only
\(\ell_0=c_{\rm int}t_0\), and
\([G]=L^3M^{-1}T^{-2}\). The dimension is fixed by the line; the
mathematical problem is to select the dimensionless character \(\theta\).

The CT108 calendar defines

\begin{equation}
\label{rm:eq:gravity-clock108}
T_{108}=108t_0,
\qquad
\omega_{108}=\frac{2\pi}{108t_0},
\qquad
R_{108}=\frac{c_{\rm int}}{\omega_{108}}
=\frac{54}{\pi}\ell_0.
\end{equation}

For each action sheet, let

\begin{equation}
\label{rm:eq:gravity-clock-energy-mass}
E_{108,a}=\hbar_a\omega_{108},
\qquad
m_{108,a}=\frac{E_{108,a}}{c_{\rm int}^2}.
\end{equation}

Then, the reduced Compton wavelength coincides with the radius of the cycle:

\begin{equation}
\label{rm:eq:gravity-compton-clock}
\bar\lambda_{C,a}
=\frac{\hbar_a}{m_{108,a}c_{\rm int}}
=R_{108}.
\end{equation}

\begin{theorem}[Uniqueness of the selector under the HMT periodicity condition]
\label{rm:thm:gravity-circular-selector}
With the reduced convention
\(r_g=Gm/c_{\rm int}^2\) and the HMT periodicity premise that identifies the gravitational
radius of the CT108 quantum with its phase radius, the closure condition
\begin{equation}
\label{rm:eq:gravity-closing-condition}
r_{g,a}(\theta)=R_{108}
\end{equation}
selects a unique positive solution,
\begin{equation}
\label{rm:eq:gravity-theta108}
\boxed{
\theta_{108}^{\rm circ}=\left(\frac{54}{\pi}\right)^2
=295.4525715010569415303525\ldots .}
\end{equation}
In that embodiment,
\begin{equation}
\label{rm:eq:gravity-four-lengths}
R_{108}=\bar\lambda_{C,a}=r_{g,a}=\ell_P^{\rm circ}.
\end{equation}
\end{theorem}

\begin{proof}
By \cref{rm:eq:gravity-transducer,rm:eq:gravity-clock-energy-mass},
\[
r_{g,a}(\theta)
=\frac{G_a(\theta)m_{108,a}}{c_{\rm int}^2}
=\theta c_{\rm int}t_0^2\omega_{108}
=\theta\frac{\pi}{54}\ell_0.
\]
Comparing it with
\(R_{108}=(54/\pi)\ell_0\) necessarily gives
\(\theta=(54/\pi)^2\). The equation is linear in \(\theta\) and all the
factors are positive, so the solution is unique.
\end{proof}

The equality \(r_g=R_{108}\) does not follow from the dimensionality of \cref{rm:eq:gravity-transducer}. It is the explicit geometric condition of this HMT realization: action, Compton, and periodic return must determine the same length. The theorem proves that, once this structural condition is accepted, the character \(\theta\) is fixed without using the SI value of \(G\).

\begin{remark}[Radius convention]
The theorem uses the reduced gravitational radius \(Gm/c^2\). If it is replaced
by the Schwarzschild radius \(2Gm/c^2\), the selector changes by a factor of
two. The two expressions correspond to distinct geometric conditions,
so the convention must be declared before the calculation.
\end{remark}

\subsection{Planck System and Bidirectional Covariance}

\begin{corollary}[Planck Family Associated with the CT108 Chronological Structure]
\label{rm:cor:gravity-planck-family}
For \(G_a=G_a(\theta_{108}^{\rm circ})\),
\begin{align}
\ell_P&=\frac{54}{\pi}\ell_0,
&t_P&=\frac{54}{\pi}t_0,
\label{rm:eq:gravity-planck-length-time}\\
E_{P,a}&=\frac{\hbar_a}{t_P}
=\frac{\pi\hbar_a}{54t_0}
=\frac{h_a}{108t_0},
&\ell_P^2&=\frac{G_a\hbar_a}{c_{\rm int}^3}
=\theta_{108}^{\rm circ}\ell_0^2.
\label{rm:eq:gravity-planck-energy-area}
\end{align}
Adem\'as,
\begin{equation}
\label{rm:eq:gravity-planck-force-power-density}
F_{P,a}=\frac{c_{\rm int}^4}{G_a},
\qquad
P_{P,a}=\frac{c_{\rm int}^5}{G_a},
\qquad
\rho_{P,a}=\frac{\hbar_a}
{(\theta_{108}^{\rm circ})^2c_{\rm int}^5t_0^4}.
\end{equation}
\end{corollary}

\begin{proof}
The first three identities follow from
\(\ell_P=R_{108}\), \(t_P=\ell_P/c_{\rm int}\), and
\(h_a=2\pi\hbar_a\). For the area,
\[
\frac{G_a\hbar_a}{c_{\rm int}^3}
=\theta_{108}^{\rm circ}c_{\rm int}^2t_0^2
=\theta_{108}^{\rm circ}\ell_0^2.
\]
The remaining expressions are the dimensional compositions of force,
power, and mass per volume with the same section.
\end{proof}

Let
\begin{equation}
\label{rm:eq:gravity-Ract}
R_{\rm act}:=\frac{\hbar_{\rm pre}}{\hbar_{\rm ret}}.
\end{equation}
Since the action appears in the numerator of the chronological mass and in the
denominator of \(G\),
\begin{equation}
\label{rm:eq:gravity-twoway-covariance}
\frac{m_{108,\rm pre}}{m_{108,\rm ret}}=R_{\rm act},
\qquad
\frac{G_{\rm pre}}{G_{\rm ret}}=R_{\rm act}^{-1},
\qquad
G_a\hbar_a
=\theta_{108}^{\rm circ}c_{\rm int}^5t_0^2.
\end{equation}
Therefore, \(\ell_P^2=G_a\hbar_a/c^3\) does not change sheets. Gravity and
action are reciprocally oriented, while area preserves the
common geometry. This cancellation is the interface of confluence toward the area
operator of the \cref{rm:ch:modular}.

\subsection{Icosahedral Component of Dimension \(5\)}

The antecedent is neither \(\R\) nor a subsequent selection of five
polarizations. The object responsible for transporting the incidence is defined
explicitly before introducing the exceptional realization.

Let \(X_{\rm TPK}^{\rm exc}\) be the subspace of the enriched state \(\TPK\)
that preserves, in addition to the arithmetic output, the orientation, route,
carry, and exceptional incidence. To make the geometric target space explicit,
let \(\varphi_g=(1+\sqrt5)/2\), \(\nu=\sqrt{1+\varphi_g^2}\), and define
\begin{equation}
\label{rm:eq:gravity-icosahedron-vertices}
\Omega_{12}^{\rm ico}
=\frac1\nu\bigl(
\{0\}\times\{\pm1\}\times\{\pm\varphi_g\}
\;\cup\;
\{\pm1\}\times\{\pm\varphi_g\}\times\{0\}
\;\cup\;
\{\pm\varphi_g\}\times\{0\}\times\{\pm1\}
\bigr)\subset S^2.
\end{equation}
Its twelve elements are the oriented vertex classes of an icosahedron.
The incidence is not left verbal: it is fixed by
\begin{equation}
\label{rm:eq:gravity-icosahedron-incidence}
\mathcal I_{12}^{\rm ico}
=\{(v,w)\in\Omega_{12}^{\rm ico}\times\Omega_{12}^{\rm ico}:
v\neq w,\ \langle v,w\rangle=1/\sqrt5\}.
\end{equation}
Each vertex has five neighbors in this relation. The reading bridge
is typed as a map
\begin{equation}
\label{rm:eq:gravity-exceptional-reader-map}
\pi_{\rm ico}:X_{\rm TPK}^{\rm exc}\longrightarrow\Omega_{12}^{\rm ico}.
\end{equation}
The cardinality twelve of the image of \(\pi_{\rm ico}\) is not sufficient.
To define an incidence realization, the mapping must satisfy
simultaneously:
\begin{enumerate}[label=(\roman*)]
\item being surjective and constant exactly on the fibers that
discard the arithmetic chart but preserve the exceptional route;
\item intertwine the action of the oriented automorphisms of the antecedent
with a transitive action on \(\Omega_{12}^{\rm ico}\);
\item transporting the declared TPK incidence to
\(\mathcal I_{12}^{\rm ico}\), not
merely its cardinalities;
\item preserve the leaf involution and the nonadic return in the
stabilizer of the corresponding fiber.
\end{enumerate}

These conditions separate three assertions: the existence of the map
\(\pi_{\rm ico}\), the identification of the group that acts, and the
linear representation constructed thereafter. The second and third
cannot be used to define the first retrospectively.

\begin{theorem}[Incidence Descent to the Icosahedral Quotient]
\label{rm:thm:gravity-exceptional-descent}
Suppose that \(\pi_{\rm ico}\) satisfies (i)--(iv) and that the oriented action
induced on \(\Omega_{12}^{\rm ico}\) has group
\(G_{\rm ico}\simeq A_5\). Then, for every class
\(v_0\in\Omega_{12}^{\rm ico}\),
\begin{equation}
\label{rm:eq:gravity-icosahedral-quotient}
H_{v_0}=\operatorname{Stab}_{G_{\rm ico}}(v_0)\simeq C_5,
\qquad
\Omega_{12}^{\rm ico}\simeq G_{\rm ico}/H_{v_0}\simeq A_5/C_5.
\end{equation}
The application
\begin{equation}
\label{rm:eq:gravity-coset-map}
\kappa_{v_0}:G_{\rm ico}/H_{v_0}\longrightarrow\Omega_{12}^{\rm ico},
\qquad gH_{v_0}\longmapsto g\cdot v_0,
\end{equation}
is an equivariant bijection and preserves incidence if, in the quotient,
one declares
\begin{equation}
\label{rm:eq:gravity-coset-incidence}
gH_{v_0}\sim g'H_{v_0}
\quad\Longleftrightarrow\quad
(g\cdot v_0,g'\cdot v_0)\in\mathcal I_{12}^{\rm ico}.
\end{equation}
Therefore, the composite \(\kappa_{v_0}^{-1}\pi_{\rm ico}\) is the typed map from the TPK
antecedent to the quotient \(A_5/C_5\).
\end{theorem}

\begin{proof}
The list of \cref{rm:eq:gravity-icosahedron-vertices} has twelve elements. By transitivity, \(|G_{\rm ico}\cdot v_0|=12\), and the orbit--stabilizer theorem gives \(|H_{v_0}|=60/12=5\). Every group of prime order five is cyclic; hence \(H_{v_0}\simeq C_5\). The formula for \(\kappa_{v_0}\) does not depend on the representative: if \(gH_{v_0}=g'H_{v_0}\), then \(g^{-1}g'\in H_{v_0}\) and \(g'v_0=gv_0\). Surjectivity follows from transitivity, and equality of the cardinalities gives injectivity. Equivariance is immediate: \(\kappa_{v_0}(agH_{v_0})=a\kappa_{v_0}(gH_{v_0})\). Finally, \cref{rm:eq:gravity-coset-incidence} is precisely the transport of structure through that bijection; hypothesis (iii) makes this transport commute with \(\pi_{\rm ico}\).
\end{proof}

The theorem does not replace the construction of \(\pi_{\rm ico}\): it proves
that it is obtained once its four requirements have been verified. In the complete
HMT profile, the incidence corridor reaching
Paley--Witt--Golay--Leech--\(V^\natural\)--Monster--Moonshine is the
domain of definition of that map. In this autonomous volume, the exact input is
the map \cref{rm:eq:gravity-exceptional-reader-map} with its conditions; the entire
segment from \(A_5/C_5\) to TT reduction is proved below
without again consulting a gravitational value.

The chain of provenance is expressed through all its morphisms:
\begin{equation}
\label{rm:eq:gravity-exceptional-provenance}
\APP\longrightarrow\TRIT\longrightarrow\TPK
\longrightarrow X_{\rm TPK}^{\rm exc}
\xrightarrow{\ \pi_{\rm ico}\ }\Omega_{12}^{\rm ico}
\xleftarrow[\simeq]{\ \kappa_{v_0}\ }A_5/C_5.
\end{equation}
The complete exceptional corridor is proved in its own dedicated development. The hypothesis that this chapter cannot replace with a count is the existence of \(\pi_{\rm ico}\) satisfying (i)--(iv). With that input fixed, descent to the quotient is \cref{rm:thm:gravity-exceptional-descent}, and \(P_5\), \(\Gamma_5\) and \(\Lambda(n)\) are constructed from it. Realizing the resulting carrier as a gravitational field remains a typed physical interface, not a consequence of cardinality.

Let \(G=A_5\), let \(H\simeq C_5\) be the stabilizer of an oriented vertex
of the icosahedron, and consider
\begin{equation}
\label{rm:eq:gravity-v12}
V_{12}=\R[G/H],
\qquad |G/H|=\frac{60}{5}=12.
\end{equation}
The permutation representation \(\rho_{12}\) decomposes as
\begin{equation}
\label{rm:eq:gravity-v12-decomposition}
V_{12}\simeq\mathbf1\oplus V_3\oplus V_{3'}\oplus V_5.
\end{equation}
The central idempotent of the character summand
\(\chi_5=(5,1,-1,0,0)\) is
\begin{equation}
\label{rm:eq:gravity-p5}
\boxed{
P_5=\frac5{60}\sum_{g\in A_5}\chi_5(g^{-1})\rho_{12}(g)
=\frac1{12}\left(
5I+\sum_{g\in2A}\rho_{12}(g)
-\sum_{g\in3A}\rho_{12}(g)
\right).}
\end{equation}

\begin{theorem}[Rank-five gravitational projector and coupler]
\label{rm:thm:gravity-p5-coupling}
The operator \(P_5\) satisfies
\begin{equation}
\label{rm:eq:gravity-p5-properties}
P_5^2=P_5=P_5^*,
\qquad \Tr P_5=\rank P_5=5.
\end{equation}
For
\begin{equation}
\label{rm:eq:gravity-G5}
\mathbb G_{5,a}:=G_a(\theta_{108}^{\rm circ})P_5
\end{equation}
we have
\begin{equation}
\label{rm:eq:gravity-G5-spectrum}
\Spec(\mathbb G_{5,a})
=\{G_a(\theta_{108}^{\rm circ})\text{ with multiplicity }5,
\;0\text{ with multiplicity }7\}.
\end{equation}
\end{theorem}

\begin{proof}
Character orthogonality converts \cref{rm:eq:gravity-p5} into the
central idempotent that acts as the identity on the unique copy of
\(V_5\) and as zero on the other three summands of
\cref{rm:eq:gravity-v12-decomposition}. Reality of the character and
orthogonality of \(\rho_{12}\) give self-adjointness. Its trace and rank
are five. Multiplying a rank-five projector by the positive scalar
\(G_a\) immediately produces the indicated spectrum.
\end{proof}

The expression \(\mathbb G_{5,a}\) does not add five different values of
the gravitational constant. It is the same section of \(G\) acting on
the selected tensor summand and annihilating the other seven modes of the
module under the permutation action.
